\subsection{Problemas de Práctica}

% --- Intro ---
\begin{frame}{Problemas de Práctica}
    Es fundamental incluir una lista de problemas para reforzar el tema. 
    El comando \texttt{\textbackslash practiceproblem} facilita este formato estandarizado.

    \vspace{0.3cm}
    \textbf{Argumentos:} \texttt{\{Nombre\}\{Dificultad\}\{Juez\}\{Link\}}
\end{frame}

% --- Lista de problemas ---
\begin{frame}{Lista de Problemas}
    \begin{itemize}
        \practiceproblem{Sum of Two Values}{Easy}{CSES}{https://cses.fi/problemset/task/1640}
        \practiceproblem{Maximum Subarray Sum}{Medium}{CSES}{https://cses.fi/problemset/task/1643}
        \practiceproblem{Counting Tilings}{Hard}{CSES}{https://cses.fi/problemset/task/2181}
    \end{itemize}
\end{frame}

\begin{frame}{Lista de Problemas}
    \begin{itemize}
        \practiceproblem{Hamiltonian Flights}{Very Hard}{CSES}{https://cses.fi/problemset/task/1690}
        \practiceproblem{Problema Personalizado}{Unknown}{Local}{}
    \end{itemize}
\end{frame}
